\documentclass[uplatex,dvipdfmx]{jsarticle}
\usepackage{amsmath,amssymb}
\usepackage{graphicx}
\usepackage{here}
\usepackage{url}
\usepackage{listings,jvlisting}
\lstset{
  basicstyle={\ttfamily\small},
  frame=tRBl,
  framesep=5pt,
  breaklines=true,
  columns=[l]{fullflexible},
  numbers=left,
  numberstyle={\scriptsize},
  numbersep=5pt,
  xrightmargin=0zw,
  xleftmargin=3zw,
  lineskip=-0.5ex
}

\title{インタラクティブメディア 最終課題 レポート}
\author{学籍番号 : 2331033\\氏名 : 桑島 大知}
\date{\today}

\begin{document}
\maketitle
\section{作品について}
今回、インタラクティブメディアの最終課題として、常に前進する魚を操作して対戦するゲームを作成した。
作品は完成せず、魚を2匹操作して遊ぶことができる状態までである。
チンアナゴを避けながら、相手の魚よりチンアナゴに当たらずに残っていると勝利するというルールである。
\section{工夫した点}
第一講時の最終課題説明のように線画で作品を作ることに注力した。

オブジェクト指向で記述し、各クラスの役割を明確にした。
また、魚を2匹操作できるようにし、対戦できるようにした。
これによりプログラムの管理をしやすくなり、拡張性を高めることで完全な感性ではなくとも遊べる状態まで持っていくことができ、
アジャイル的開発を実現し当初の構想まで持っていくことができなくとも問題なく動くようにできた。
また、再利用性を高めることができ、文字列などの表示に関するクラスを汎用的に作成できた。

回転は四元数を用いて実装した。一方で透視投影変換、カメラ座標変換は同時座標系の行列で実装することになった。

魚の加速機能で単純にの速度を定数倍するのではなく、上限をもうけ加減速度的に変化させるようにしたことで、より自然な動きになった。
また、線画では人間の認知的に上下前後関係が分かりずらいため、魚の影として体の軸と曲がり方を床面に投影して描画することで、魚の位置関係が分かりやすくなるように工夫した。
全体的に線画で3Dを扱うためのインタラクションを意識した作品になった。
\section{苦労した点}
多くのことを妥協することになった。まず魚の動きであるが、魚特有の動きを実装することができなかった。本来は数学的に正しい魚の体の動かし方を調査し実装することが目的の一つであったが、時間が足りず、単純にキーボード入力に応じて魚の向きを変え、直進するだけの動きになってしまった。
ただし、視覚的に魚がどううごいているのかを分かりやすくするために、魚の体を構成する頂点をサイン波で変位させることで、魚が泳いでいるように見えるようにした。
また、魚を2匹操作できるようにしたが、入力キーを変えるためにメソッドを複製する必要があり、冗長なコードになってしまった。
そして、チンアナゴも描画する際に円柱で表現する簡易的な表現にとどまったが、視覚的に動いているのが分かりやすいようにx-z平面に描いた円を毎フレーム再計算して線形に移動させて描画した。

今まで作成したコードのリファクタリングに時間がかかり、追加機能の実装に時間をたくさん割くことができなかった。
\section{実装}
\subsection{操作方法}
1PはWASDキーで上下左右を操作し、スペースキーで加速するようにした。
2Pは矢印キーで上下左右を操作し、右コントロールキーで加速するようにした。
\subsection{クラス構成}
\begin{itemize}
\item GameManager: ゲーム全体を管理するクラス。各オブジェクトの生成、更新、描画を行う。
\item Fish: 魚を表現するクラス。位置、速度、向きの管理と更新、描画を行う。
\item ChineseEel: チンアナゴを表現するクラス。位置、速度の管理と更新、描画を行う。
\item Frame: 水槽の枠を表現するクラス。位置の管理と描画を行う。
\item Camera: カメラを表現するクラス。視点の管理と更新を行う。
\item StringsBlock: 文字列表示を行うクラス。位置、フォントサイズ、色の管理と描画を行う。
\end{itemize}
また、補助関数のファイルとして、四元数の計算を行う複数の関数をもつ\verb|quarternions.py|、簡易viewport clipping計算する関数、x-z平面に円を描く関数を持つ\verb|object_module.py|、ウィンドウ管理を行い実行するmain関数をもつ実行部分である\verb|main.py|の4つを実装した。

\subsection{コード}
まず、実行部分である\verb|main.py|のコードを示す。このファイルの77行目以降が実行部分である。
左上を原点としx座標が右方向に、y座標が下方向に増加するものであった座標系を中心原点のx軸が右方向に、y軸が上方向に増加する座標系に変換している課題用に配布されたコードを流用している。
\lstinputlisting[language=Python,caption=main.py,label=main]{main.py}

次に魚を表現するクラスである\verb|Fish.py|のコードを示す。コンストラクタは23行目からである。
コンストラクタでは魚の初期位置の中心三次元座標、速度三次元ベクトル、色のタプル、魚の大きさ、骨の分割数、四元数で初期向きを受け取り、各種パラメータを初期化している。
ここで受け取った三次元座標を四元数として保持する。速度ベクトルについても同様である。 魚の向きは四元数で保持している。
また、加速したときの最大速度、最低速度を初速度から計算して保持している。
メンバ変数であるdecoreationは魚の体を構成する頂点群を繋げて描画する単位ごとにリストにして、リストのリストとして保持している。これは体の長さや骨の分割数に応じて計算されている。

39行目のset\_decorationメソッドでは骨の一本の線しかない魚に対して、体を構成する頂点群を計算して魚らしい外見にするメソッドである。

70行目から魚を描画するdisplayメソッドが定義されている。引数として描画先のバッファ、ビュー変換行列、透視投影変換行列、画面の座標に変換するスケール、原点座標を受け取る。
描画にはまず魚のモデリング座標と四元数の姿勢から回転をし、四元数の4次元座標ベクトルから同次座標系の4次元座標ベクトルに変換し、PPMとビュー行列を前からかけて描画する座標に変換し、簡易viewport clippingを行い、画面の座標に変換してから線画を描画している。
その後、魚の上下をわかりやすくするとさかのようなものを表示するメソッドdisplay\_landmark、身体の装飾を描画するdisplay\_decorationメソッド、魚の現在位置が分かりやすくなるように床面に影を投影して描画するdisplay\_shadowメソッドを呼び出している。

155行目から魚の状態を更新するupdateメソッドが定義されている。魚が前進していて何らかのキー入力がある場合、inputメソッドでキー入力処理を行う。
速度に回転を状態のぶんだけ加え、オイラー法で更新している。update2Pメソッドについては入力処理のinputメソッドの代わりにinput2Pメソッドを呼び出す以外は同じである。

184行目のinputメソッドでは1Pのキー入力処理を行っている。WASDキーで上下左右の回転を行い、スペースキーで加速する。
input2Pメソッドでは矢印キーで上下左右の回転を行い、右コントロールキーで加速する。
なお、updateとinputはともに入力としてキー入力、キーで外から干渉できるときの速度、オイラー法における時間刻み幅を受け取る。

death\_judgeメソッドでは魚が入力のポジションという円柱の軸のてっぺんに点を取り、そのxz平面方向に円形を描き、軸からの距離を計算し、接触判定を行うメソッド。xz平面におけるy軸平行のチンアナゴの中心軸からの距離がチンアナゴの半径より小さい場合、接触したと判定しTrueを返す。

\lstinputlisting[language=Python,caption=Fish.py,label=fish]{Fish.py}

次に、チンアナゴを表現するクラスである\verb|ChineseEel.py|のコードを示す。

コンストラクタは7行目からである。
コンストラクタではチンアナゴの初期位置の中心三次元座標、速度三次元ベクトル、色のタプル、チンアナゴの大きさを受け取り、各種パラメータを初期化している。
チンアナゴの向きはy軸方向に固定されているため、四元数で保持していない。なお、大きさは半径と最長時の長さで指定できる。

15行目からdisplayメソッドが定義されている。引数として描画先のバッファ、ビュー変換行列、透視投影変換行列、画面の座標に変換するスケール、原点座標を受け取る。
チンアナゴは地面から伸びてくるため、ちゃんと地表に出たもののみ描画したい。ポリゴンならば陰面消去やZバッファで実現できるが、今回ポリゴンのように辺を決まった順で定義していないので、難しい。
そこで、チンアナゴの先端のy座標がどこにあるかを判定し、地面より下にある部分を計算的に描画しないようにした。
その後、チンアナゴのモデリング座標を同次座標系の4次元座標ベクトルに変換し、PPMとビュー行列を前からかけて描画する座標に変換し、簡易viewport clippingを行い、画面の座標に変換してから線画を描画している。

102行目からチンアナゴの状態を更新するupdateメソッドが定義されている。地面から伸びてきて、指定した最大長さに達していない場合、y座標方向に速度を加え、オイラー法で更新している。
また、チンアナゴが伸び切った後は速度反転し、地面に戻るようにしている。なお、魚の死亡のためにget\_positionメソッドで現在位置を取得できるようにしている。

\lstinputlisting[language=Python,caption=ChineseEel.py,label=chineseeel]{ChineseEel.py}

次にFrameクラスを示す。水槽の枠を表現するクラスである。
コンストラクタは8行目からである。
コンストラクタでは枠の中心三次元座標、フレームの姿勢を表す四元数、垂直方向つまり奥行き方向の分割数、水平つまり横方向の分割数、垂直方向の全体の大きさ、水平方向の全体の大きさを受け取り、各種パラメータを初期化している。
フレームの姿勢を表す四元数は今回初期値での変更以外で変えようがなく、斜面でもできるようにしたかったため設定した。
垂直方向つまり奥行き方向の分割数、水平つまり横方向の分割数、垂直方向の全体の大きさ、水平方向の全体の大きさは保持せず、
直接描画する辺の頂点群のモデリング座標に変換してリストのリストとして保持している。

37行目からdisplayメソッドが定義されている。引数として描画先のバッファ、ビュー変換行列、透視投影変換行列、画面の座標に変換するスケール、原点座標、色のタプルを受け取る。
他のオブジェクトと違い色は引数で受け取るようにしている。

\lstinputlisting[language=Python,caption=Frame.py,label=frame]{Frame.py}

Cameraクラスを示す。カメラを表現するクラスである。
コンストラクタは8行目からである。
現在位置、注視点、視界のbottom,left,top,right,nearの7つのパラメータを受け取り、各種パラメータを初期化している。
このパラメタから透視投影変換行列を計算し、保持している。
また、ビュー変換行列も計算し、保持している。

inputメソッドなどが定義されている。本来視点を動かせるようにしていたが、インタラクションを考えて全体カメラ一つに固定した。
よって、実際今回は変換行列を計算して保持、出力するだけのクラスになっている。

\lstinputlisting[language=Python,caption={Camera.py},label=Camera]{Camera.py}

StringsBlockクラスを示す。文字列表示を行うクラスである。
コンストラクタは13行目からである。
文字列、フォントサイズ、二次元座標、色のタプルを受け取り、各種パラメータを初期化している。
コンストラクタではそれ以降のAからZまでの文字の頂点を計算するメソッドから頂点群を計算し、リストを作成して保持している。

28行目からdisplayメソッドが定義されている。引数として描画先のバッファ、画面の座標に変換するスケール、原点座標を受け取る。
リストの要素である文字単位の頂点群を順に取り出し、print\_stringメソッドで画面に描画している。

36行目からprint\_stringメソッドが定義されている。引数として描画先のバッファ、画面の座標に変換するスケール、原点座標、文字単位の頂点群を受け取る。
各頂点群を画面の座標に変換し、線画を描画している。

\lstinputlisting[language=Python,caption={StringsBlock.py},label=StringsBlock]{StringsBlock.py}

次にGameManagerクラスを示す。ゲーム全体を管理するクラスである。
12行目からコンストラクタが定義されている。
コンストラクタでは魚2匹、チンアナゴ複数、水槽の枠、カメラ、文字列表示オブジェクトを生成し、時間刻み幅dtなど各種パラメータを初期化している。

34行目からupdateメソッドが定義されている。引数としてキー入力、時間刻み幅を受け取る。
このメソッドで各オブジェクトの状態更新、描画、判定を行っていて、いわばゲームループの本体である。
まず、魚2匹、チンアナゴ複数、水槽の枠、カメラの状態をupdateしている。魚を操作しているときに難しくなってほしいので、キー入力があるときチンアナゴが増えるように、
カウントダウンを行い、一定時間ごとにチンアナゴを追加している。

チンアナゴと魚からの接触判定を行い、接触した魚を死亡扱いにしている。
最後に文字列表示オブジェクトの内容を更新している。

基本的にゲームに追加機能を実装するときはそれぞれオブジェクトを作ったのち、このクラスのupdateメソッドに追加していくことになる。

\lstinputlisting[language=Python,caption={GameManager.py},label=GameManager]{GameManager.py}

最後に、補助関数群を示す。quaternions.py、object\_module.pyである。それぞれ、簡潔な説明をコメントアウトとしてつけている。
まず、quaternions.pyのコードを示す。
\lstinputlisting[language=Python,caption={quaternions.py},label=quaternions]{quaternions.py}

次にobject\_module.pyのコードを示す。
\lstinputlisting[language=Python,caption={object\_module.py},label=objectmodule]{object_module.py}

以上で、コードの説明を終わる。


\end{document}